\scanpage{29}

% head, left
$\Gamma=\mathbb G$

% head, right
$M\times B'\qquad \operatorname{Ext}(B',DM)$

% head, second line
$\deleted{\operatorname{Hom}(B',\operatorname{Ext}(\unsure{M|B},\unsure{G|A|\mathbb G_m}))}=0$

% upper left and centre
\diagram{\begin{tikzcd}[ampersand replacement=\&,column sep=1.6em,row sep=1.7em]
B \& B' \& \text{une biext. }\inserted{\mathcal E}\text{ vers }\mathbb G_m\\
M\arrow[u,"i"] \& M'\arrow[u,"i'"'] \&
\end{tikzcd}}

% upper pair of maps
\[
\begin{aligned}
 B'&\xrightarrow{\alpha'}\underline{\operatorname{Ext}}(B,G_m)\mathbin{-}\underline{\operatorname{Ext}}(M,G_m)\quad\text{nul}\\
 B&\mathrel{\overset{\alpha}{\text{---}}}\underline{\operatorname{Ext}}(B',G_m)\mathbin{-}\underline{\operatorname{Ext}}(M',G_m)\quad\text{nul}
\end{aligned}
\]

% right of upper maps
\margin{$\begin{array}{l}\text{p. ex }\operatorname{Hom}(B',\underline{\operatorname{Ext}}(M,\overset{\substack{0\\\Vert}}{G_m}\\ \text{p.ex }\operatorname{Hom}(B,\underline{\operatorname{Ext}}(M',\underset{\substack{\Vert\\0}}{G_m}\end{array}$}

% upper middle
d'où extensions

% two extension rows
\diagram{\begin{tikzcd}[ampersand replacement=\&,column sep=1em,row sep=.35em]
\cdot\arrow[r,dash] \& DM\arrow[d,equal]\arrow[r,"j'"] \& G'\arrow[r,"p'",dash] \& B'\arrow[r] \& 0\\
\& T' \& \& \&\\
0\arrow[r] \& DM'\arrow[d,equal]\arrow[r,"j",dash] \& G\arrow[r,"p",dash] \& B\arrow[r] \&0\\
\&T\&\&\&
\end{tikzcd}}

% right of extensions, upper
\margin{$\mathcal E_{M,G'}$ est canoniqt splitté par $\rho$}

% right of extensions, lower
\margin{$\mathcal E_{G,M'}$ est can. splitté $\rho'$}

% example
\uline{Exemple}\qquad Si \deleted{\unsure{$B$ est}} $\underline{\alpha'}$ \deleted{\unsure{éta}}\inserted{est} un isom, alors $i'$ \unsure{revient à donner|revient : la donnée|revient au même de donner} un $i'\in\operatorname{Hom}(M',\underline{\operatorname{Ext}}(B,G_m))$

% example, arrows and conditions
\diagram{\begin{tikzcd}[ampersand replacement=\&,column sep=.5em,row sep=.45em]
{}\& \arrow[l,shift left=.4ex]\arrow[r,shift left=.4ex] \& \operatorname{Biext}(B,M';G_m) \& \arrow[l,shift left=.4ex]\arrow[r,shift left=.4ex] \& \deleted{\unsure{\operatorname{repr}}}\inserted{\operatorname{Ext}(B,D(M'))}\\
\&\text{si }\deleted{\unsure{\operatorname{Hom}}}\operatorname{Ext}(M',\underline{\operatorname{Hom}}(B,G_m))=0 \&\& \text{si }\operatorname{Hom}(B,\underline{\operatorname{Ext}}(M',G_m))=0 \&\\
\&\text{p.ex }\underline{\operatorname{Hom}}(B,G_m)=0\&\&\text{p.ex. }\underline{\operatorname{Ext}}(M',G_m)=0\&
\end{tikzcd}}

% middle, continuation
aussi \unsure{bien|pour|donc}. Donc la donnée de $i'$ revient à celle de l'ext. $G$. Et réciproquement

% lower half
Considérons l'image inverse \inserted{$\mathcal E_{\unsure{i,i'|i',i'}}$} de $\mathcal E$ sur $(M,M')$. Alors des splittings \inserted{$\rho_i$} de cette biextension correspondent aux hom $j$ de $M$ dans $G$ relevant $i$, ou aux \unsure{mor}$\,j'$ de $M'$ dans $G'$ relev. de $i'$.

% interlinear diagram
\diagram{\begin{tikzcd}[ampersand replacement=\&,column sep=1.4em,row sep=.2em]
 M\arrow[r,dash]\& G\arrow[r,dash]\& B\\
 M'\arrow[r,dash]\&G'\arrow[r,dash]\&B'
\end{tikzcd}}

% above cancellation
\deleted{\unsure{sur le}}

% lower text and rewritten block
Des relèvements \deleted{\unsure{Si}} $j$ et $j'$ \deleted{\unsure{se}} correspondent ainsi, \deleted{alors \unsure{on trouve un} splitting canonique \deleted{\unsure{de}} $\sigma$ de $\mathcal E_{G,M'}$ et un splitting canonique $\sigma'$ de $\mathcal E_{M,G'}$, ces deux splittings coïncident sur $M,M'$ \unsure{avec le} splitting \unsure{donné|associé} $h$ de $\mathcal E_{M,M'}$.}

% last three lines
ssi les deux splittings de $\mathcal E_{M,M'}$, déduits d'une part du splitting $\rho$ de $\mathcal E_{M,G'}$ et de $j':M'\mathbin{-}G'$, d'autre part du splitting $\rho'$ de $\mathcal E_{G,M'}$ et de $j:M\mathbin{-}G$, coïncident.
% archive sheet number: 28
